\section{Morris Inorder Traversal}
\label{sec:trees-morris-inorder}

\problemstatement{Return the inorder traversal of a binary tree using
$O(1)$ extra space -- no recursion and no explicit stack.}

\begin{patternbox}
\emph{``$O(1)$ space,'' ``no stack,'' ``no recursion''} points to
\textbf{Morris traversal}, which recovers the ``go back up to the
ancestor'' capability -- normally provided by a stack -- by temporarily
rewiring the tree's own unused null pointers into \textbf{threads} back to
the ancestor, then undoing them.
\end{patternbox}

\subsection*{Understanding the problem carefully}

The only reason inorder needs a stack is to remember, after finishing a
left subtree, which node to return to. Morris's insight: the node we must
return to is exactly the \emph{inorder predecessor}'s successor -- and the
inorder predecessor of a node (the rightmost node of its left subtree) has
a null right pointer we can borrow. We point that null back at the current
node, creating a temporary ``thread.'' After the left subtree is walked, we
arrive back via the thread, record the node, remove the thread (restoring
the tree), and move right.

\begin{ideabox}[Thread the Predecessor's Null Right Pointer Back to You]
At each node with a left child, find its inorder predecessor (go left once,
then all the way right). If the predecessor's right is null, set it to the
current node (make the thread) and move left. If it is already the current
node, the left subtree is done: remove the thread, record the node, and
move right. Nodes without a left child are recorded immediately and you
move right.
\end{ideabox}

\subsection*{Algorithm}

\begin{enumerate}
  \item \textit{curr} $=$ root.
  \item While \textit{curr} is non-null:
  \begin{enumerate}
    \item If \textit{curr} has no left child: record \textit{curr}; move
          \textit{curr} $=$ \textit{curr}.right.
    \item Else find \textit{pred} $=$ rightmost node of
          \textit{curr}.left (not passing through an existing thread):
    \begin{enumerate}
      \item If \textit{pred}.right is null: set it to \textit{curr};
            move \textit{curr} $=$ \textit{curr}.left.
      \item Else (\textit{pred}.right $=$ \textit{curr}): reset
            \textit{pred}.right to null; record \textit{curr}; move
            \textit{curr} $=$ \textit{curr}.right.
    \end{enumerate}
  \end{enumerate}
\end{enumerate}

\subsection*{Dry run}

\dryrun{Tree: root $2$; left $1$; right $3$.}

\begin{itemize}
  \item \textit{curr}$=2$, has left $1$. Predecessor of $2$ is $1$;
        $1$.right null $\to$ thread $1\!\to\!2$; move left to $1$.
  \item \textit{curr}$=1$, no left $\to$ record $1$; move right via
        thread to $2$.
  \item \textit{curr}$=2$, has left; predecessor $1$.right $=2$ (thread)
        $\to$ remove thread, record $2$, move right to $3$.
  \item \textit{curr}$=3$, no left $\to$ record $3$.
  \item Inorder: $1,2,3$.
\end{itemize}

\subsection*{C++ implementation}

\begin{lstlisting}
#include <bits/stdc++.h>
using namespace std;

struct TreeNode {
    int val;
    TreeNode* left;
    TreeNode* right;
    TreeNode(int v) : val(v), left(nullptr), right(nullptr) {}
};

vector<int> morrisInorder(TreeNode* root) {
    vector<int> out;
    TreeNode* curr = root;
    while (curr != nullptr) {
        if (curr->left == nullptr) {
            out.push_back(curr->val);        // no left: record, go right
            curr = curr->right;
        } else {
            TreeNode* pred = curr->left;
            while (pred->right != nullptr && pred->right != curr)
                pred = pred->right;          // find inorder predecessor
            if (pred->right == nullptr) {
                pred->right = curr;          // create thread
                curr = curr->left;
            } else {
                pred->right = nullptr;       // remove thread
                out.push_back(curr->val);    // left subtree done: record
                curr = curr->right;
            }
        }
    }
    return out;
}

int main() {
    TreeNode* root = new TreeNode(2);
    root->left = new TreeNode(1);
    root->right = new TreeNode(3);

    for (int x : morrisInorder(root)) cout << x << " ";
    cout << "\n";
    return 0;
}
\end{lstlisting}

\begin{complexitybox}
\textbf{Time Complexity.} $O(n)$ -- although finding predecessors seems to
add work, each edge is traversed at most a constant number of times, so
the total remains linear. \textbf{Space Complexity.} $O(1)$ -- no stack or
recursion; only a few pointers, with the tree restored to its original
shape at the end.
\end{complexitybox}

\begin{takeawaybox}
Morris trades the stack for temporary threads woven from the tree's own
null pointers, achieving true $O(1)$ space. The discipline that keeps it
correct is always removing a thread the second time you meet it, so the
tree is left exactly as it began.
\end{takeawaybox}
